\chapter{Repository and Contribution Evidence}
The repository holds the simulation sources, the extraction and training pipeline,
the automated test suite and this report. Continuous integration evidence is reported
in Section~\ref{sec:repository} and shown in Figure~\ref{fig:ci}.

\textbf{Repository.} \url{https://github.com/AdeenaZaib/VROOM}

\section*{Per-member commit counts}
Output of \texttt{git shortlog -sne}:

\begin{verbatim}
     4  AdeenaZaib <adeenazaib42@gmail.com>
     3  khaulabilqees <khbilq16@gmail.com>
\end{verbatim}

\section*{Commit log}
Output of \texttt{git log --pretty='\%h \%an \%ad \%s' --date=short}:

\begin{verbatim}
f009ac6  khaulabilqees  2026-10-13  update
31cb803  khaulabilqees  2026-10-13  update
2168990  AdeenaZaib     2026-10-13  add CI tests for feature extraction
                                    and training pipeline
856beba  AdeenaZaib     2026-10-13  adding the feature extraction stuff
ed4b0f6  AdeenaZaib     2026-10-13  added gitHub actions test workflow
8b765f6  khaulabilqees  2026-10-13  topology in omnett
0506e63  AdeenaZaib     2026-09-08  poc
\end{verbatim}

The mapping from these commits to project milestones is given in
Table~\ref{tab:milestones}, and the per-member attribution in Section~8.4.

\textbf{Note on commit granularity.} The history is coarse: seven commits, of which
two carry the message \texttt{update}. This reflects the repository having been
introduced partway through Iteration~1 rather than at its start, so earlier work was
committed in bulk once version control was in place. Commit practice for FYP-2 is to
commit per logical change with a message naming the module, in the form
\texttt{M3: confine attack labels to the target zone}, so that the milestone mapping
above can be produced mechanically rather than reconstructed.
